\documentclass[tikz,border=6pt]{standalone}
\usepackage{fontspec}
\usepackage[mathrm=sym]{unicode-math}
\setmainfont{texgyrepagella-regular.otf}[
    BoldFont = texgyrepagella-bold.otf,
    ItalicFont = texgyrepagella-italic.otf,
    BoldItalicFont = texgyrepagella-bolditalic.otf
]
\setmathfont{texgyrepagella-math.otf}[AutoFakeBold = 2.5]
\usepackage{tikz}
\usepackage{pgfplots}
\pgfplotsset{compat=1.18}
\usetikzlibrary{arrows.meta, calc, patterns}

\begin{document}
\begin{tikzpicture}[>=Stealth, line cap=round, line join=round]

  \begin{axis}[
    width=10.5cm,
    height=8.2cm,
    xmin=0, xmax=13.5,
    ymin=-2.4, ymax=9.4,
    axis x line=middle,
    axis y line=left,
    axis line style={-Stealth, line width=0.85pt},
    xtick={2, 4, 6, 8, 10, 12},
    xticklabels={, $4$, $6$, $8$, $10$, $12$},
    xticklabel style={below=2pt, font=\small},
    ytick={-2, -1, 0, 1, 2, 3, 4, 5, 6, 7, 8},
    yticklabels={$-2$, , $0$, , $2$, , $4$, , $6$, , $8$},
    yticklabel style={left=2pt, font=\small},
    tick align=inside,
    tick style={line width=0.6pt, color=black},
    major tick length=4pt,
    clip=false
  ]

    % Axis titles
    \node[anchor=north west, font=\small] at (axis cs:12.7, -0.05) {$r\,(\mathrm{fm})$};
    \node[anchor=south, font=\small, align=center] at (axis cs:0, 9.4) {$\rho^{3^-}(r)\cdot 10^3\,(e\,\mathrm{fm}^{-3})$};

    % --- EXPERIMENTAL REGION (HATCHED BAND) ---
    % Hatching fill
    \fill[pattern=north east lines, pattern color=black!75]
      (axis cs:0, 0)
      to[out=60, in=170] (axis cs:0.9, 0.72)
      to[out=-10, in=135] (axis cs:1.75, 0.0)
      to[out=-45, in=180] (axis cs:2.35, -1.35)
      to[out=0, in=-120] (axis cs:3.0, 0.0)
      to[out=60, in=170] (axis cs:3.75, 1.95)
      to[out=-10, in=180] (axis cs:4.6, 1.08)
      to[out=0, in=-115] (axis cs:5.5, 2.5)
      to[out=65, in=195] (axis cs:6.7, 7.85)
      to[out=-15, in=115] (axis cs:7.8, 3.8)
      to[out=-65, in=160] (axis cs:8.6, 0.6)
      to[out=-20, in=180] (axis cs:9.2, 0.22)
      to[out=0, in=170] (axis cs:10.5, 0.0)
      to[out=-10, in=180] (axis cs:11.3, -0.15)
      to[out=0, in=180] (axis cs:12.2, 0.0)
      -- (axis cs:12.2, 0.0)
      to[out=180, in=0] (axis cs:11.4, -0.22)
      to[out=180, in=-5] (axis cs:10.5, -0.05)
      to[out=175, in=0] (axis cs:9.2, 0.12)
      to[out=180, in=-30] (axis cs:8.55, 0.35)
      to[out=150, in=-70] (axis cs:7.75, 3.5)
      to[out=110, in=-15] (axis cs:6.7, 7.55)
      to[out=195, in=65] (axis cs:5.55, 2.2)
      to[out=-115, in=0] (axis cs:4.65, 0.68)
      to[out=180, in=0] (axis cs:3.7, 1.30)
      to[out=180, in=70] (axis cs:3.05, 0.0)
      to[out=-110, in=0] (axis cs:2.15, -1.88)
      to[out=180, in=-40] (axis cs:0.8, -0.45)
      to[out=140, in=-20] (axis cs:0, 0)
      -- cycle;

    % Draw solid boundary lines of experimental band
    \draw[line width=0.85pt]
      (axis cs:0, 0)
      to[out=60, in=170] (axis cs:0.9, 0.72)
      to[out=-10, in=135] (axis cs:1.75, 0.0)
      to[out=-45, in=180] (axis cs:2.35, -1.35)
      to[out=0, in=-120] (axis cs:3.0, 0.0)
      to[out=60, in=170] (axis cs:3.75, 1.95)
      to[out=-10, in=180] (axis cs:4.6, 1.08)
      to[out=0, in=-115] (axis cs:5.5, 2.5)
      to[out=65, in=195] (axis cs:6.7, 7.85)
      to[out=-15, in=115] (axis cs:7.8, 3.8)
      to[out=-65, in=160] (axis cs:8.6, 0.6)
      to[out=-20, in=180] (axis cs:9.2, 0.22)
      to[out=0, in=170] (axis cs:10.5, 0.0)
      to[out=-10, in=180] (axis cs:11.3, -0.15)
      to[out=0, in=180] (axis cs:12.2, 0.0);

    \draw[line width=0.85pt]
      (axis cs:0, 0)
      to[out=-20, in=140] (axis cs:0.8, -0.45)
      to[out=-40, in=180] (axis cs:2.15, -1.88)
      to[out=0, in=-110] (axis cs:3.05, 0.0)
      to[out=70, in=180] (axis cs:3.7, 1.30)
      to[out=0, in=180] (axis cs:4.65, 0.68)
      to[out=0, in=-115] (axis cs:5.55, 2.2)
      to[out=65, in=195] (axis cs:6.7, 7.55)
      to[out=-15, in=110] (axis cs:7.75, 3.5)
      to[out=-70, in=150] (axis cs:8.55, 0.35)
      to[out=-30, in=180] (axis cs:9.2, 0.12)
      to[out=0, in=175] (axis cs:10.5, -0.05)
      to[out=-5, in=180] (axis cs:11.4, -0.22)
      to[out=0, in=180] (axis cs:12.2, 0.0);

    % --- THEORY (RING + SPETH) (DASHED LINE) ---
    \draw[dashed, dash pattern=on 4.5pt off 3pt, line width=0.95pt]
      (axis cs:0, 0)
      to[out=25, in=160] (axis cs:1.0, 0.25)
      to[out=-20, in=140] (axis cs:2.1, 0.2)
      to[out=-40, in=-150] (axis cs:2.8, 0.8)
      to[out=30, in=180] (axis cs:3.8, 3.35)
      to[out=0, in=180] (axis cs:4.7, 2.65)
      to[out=0, in=-115] (axis cs:5.6, 3.8)
      to[out=65, in=195] (axis cs:6.6, 7.3)
      to[out=-15, in=115] (axis cs:7.6, 3.8)
      to[out=-65, in=160] (axis cs:8.4, 1.2)
      to[out=-20, in=180] (axis cs:9.5, 0.1)
      to[out=0, in=180] (axis cs:11.0, 0.0);

    % --- BOHR - TASSIE - MODEL (DOTTED LINE) ---
    \draw[dotted, line width=1.3pt, dash pattern=on 1.5pt off 2.5pt]
      (axis cs:2.0, 0.0)
      to[out=0, in=180] (axis cs:3.5, 0.05)
      to[out=0, in=-120] (axis cs:4.8, 0.95)
      to[out=60, in=-115] (axis cs:5.9, 3.2)
      to[out=65, in=195] (axis cs:7.2, 7.95)
      to[out=-15, in=110] (axis cs:8.3, 4.5)
      to[out=-70, in=135] (axis cs:9.2, 1.8)
      to[out=-45, in=160] (axis cs:10.2, 0.45)
      to[out=-20, in=180] (axis cs:11.8, 0.0);

    % --- LEGEND ---
    \begin{scope}[shift={(axis cs:7.2, 6.8)}]
      % Background mask
      \fill[white, opacity=0.92] (0, 0.1) rectangle (5.7, 2.4);

      % Dashed: Theory
      \draw[dashed, dash pattern=on 4.5pt off 3pt, line width=0.95pt] (0.2, 2.1) -- (1.1, 2.1);
      \node[anchor=west, font=\footnotesize] at (1.2, 2.1) {Theory (Ring\,+\,Speth)};

      % Dotted: Bohr - Tassie
      \draw[dotted, line width=1.3pt, dash pattern=on 1.5pt off 2.5pt] (0.2, 1.45) -- (1.1, 1.45);
      \node[anchor=west, font=\footnotesize] at (1.2, 1.45) {Bohr\,--\,Tassie\,--\,Model};

      % Solid: Experiment
      \draw[line width=0.85pt] (0.2, 0.7) -- (1.1, 0.7);
      \node[anchor=west, font=\footnotesize, align=left] at (1.2, 0.6) {Experiment (Yale, Mainz,\\[-1pt]Sendai, Stanford)};
    \end{scope}

  \end{axis}
\end{tikzpicture}
\end{document}
